        \section{为什么三条线被放到同一主线}
        钨股本身可以独立成立一条小金属资源线，氟化工也可以独立成立一条制冷剂配额线，工业气体则长期有电子特气国产替代线。近期市场把三者放在同一张图里，核心原因是WF6把三种能力合并成一个前道材料瓶颈：没有高纯钨粉和氟化反应能力，就没有可认证的高纯WF6；没有晶圆厂认证，名义产能也不能变成可持续利润。

        \section{技术链条}
        WF6在CVD中提供钨源，适合高深宽比结构的钨沉积或填充。先进逻辑、DRAM、3D NAND和HBM对层数、互连密度和良率提出更高要求，电子特气纯度、金属杂质、包装容器和连续供应能力成为认证门槛。报告因此把“资源价格上涨”和“半导体认证放量”分开定价。

        \begin{exhibitbox}[Mermaid产业链图的语义摘要]
        本案例的产业链图使用Mermaid语法绘制。链条从钨矿/APT和萤石/HF开始，经高纯钨粉、含氟反应和纯化进入WF6/NF3等电子特气，再进入CVD钨沉积、刻蚀清洗，最终服务先进逻辑、DRAM、3D NAND、HBM和AI算力硬件。
        \end{exhibitbox}


\begin{exhibitbox}[制程/材料功能证据]
\scriptsize
\begin{tabularx}{\exhibitboxwidth}{L{1.2cm} X X X X}
\toprule
\textbf{制程层} & \textbf{证据点} & \textbf{半导体意义} & \textbf{需求传导} & \textbf{估值用途} \\
\midrule
前道钨源 & WF6用于低压或等离子增强CVD沉积钨、钨硅化物和钨氮化物 & 证明WF6不是泛工业气体，而是前道薄膜/金属化工艺气体 & 先进逻辑、DRAM、3D NAND、HBM相关晶圆制造和薄膜工艺 & 支持产业链交汇点和直接性评级，不直接进入单公司收入或每股收益 \\
钨连接/金属化 & CVD/ALD用于制造微小钨连接和薄阻挡层，高保形金属膜用于先进钨金属化 & 钨连接、互连、晶体管和先进存储是设备商公开列示的应用场景 & 先进存储、先进封装、互连和晶体管应用 & 支持下游需求锚和工艺链位置，不替代公司订单 \\
接触孔/通孔填充 & 钨CVD用于DRAM和逻辑芯片关键接触孔/通孔结构无空洞填充 & 高深宽比填充、接触电阻和良率是钨沉积工艺的核心约束 & DRAM、逻辑芯片、0.18微米及以下接触/通孔结构，先进节点延伸关注接触电阻 & 解释先进制程为什么会放大高纯WF6供应可靠性权重 \\
供应规格/纯度控制 & WF6供应需要镍系统、专有蒸馏去除杂质、非反应性镍钢瓶和痕量金属分析 & 纯度、痕量金属、容器和供应安全是电子特气认证的一部分 & 半导体制造客户，尤其存储芯片制造商的供应安全需求 & 支持把包装、纯化、供应稳定性列入客户认证门槛 \\
\bottomrule
\end{tabularx}
\sourcenote{Linde Electronics、Lam Research、Applied Materials、Air Products公开资料整理；制程证据证明材料功能和良率位置，不替代A股公司客户认证、订单、合同平均售价、收入确认比例或单品毛利。}
\end{exhibitbox}

这张工艺表把主题叙事压实到三个可验证层级：第一，WF6是低压或等离子增强CVD中的钨源，服务钨、钨硅化物和钨氮化物薄膜；第二，设备商把钨连接、接触孔、通孔和高保形金属膜列为先进存储、互连、晶体管和逻辑芯片应用；第三，全球气体供应商强调纯化、镍系统、非反应性钢瓶和痕量金属控制。这能证明“为什么要研究WF6”，但不能证明“哪家A股公司已经拿到订单”。因此本报告把技术证据放在产业链直接性里，把订单、平均售价和单品毛利继续锁在增长业绩门禁里。
